% ==============================================================================
% Section 9: Conclusion
% ==============================================================================
\section{Conclusion}
\label{sec:conclusion}

This paper proposes frequent losers---firms that systematically participate in and lose public procurement tenders---as a novel screening marker for cover bidding in cartel arrangements. Using comprehensive data from Brazil's BEC platform (2009--2019), comprising 4.5 million tender-items and 40 million bids, we establish that FL presence is associated with significantly higher winning prices (6--9\%), inflated bidder counts, and bid patterns consistent with complementary cover bidding (Regime~1). Bajari-Ye tests confirm that FL bids violate exchangeability and conditional independence with substantially greater magnitude than non-FL bids. External validation against CADE cartel convictions shows that 193 FL firms co-participate in tenders with known cartelists at a rate 3.5 times higher than non-FL firms.

We are transparent about the limits of our identification. The FL-price association is robust to matching, sensitivity analysis, and threshold variation, but we cannot definitively rule out that FL firms select into markets that would have had high prices regardless. The staggered DiD analysis, which could in principle establish causality, does not yield clean pre-trends in our setting. We therefore frame the FL screen as a \textit{screening marker}---a firm-level pattern strongly consistent with cover bidding and validated against known cartels---rather than as a causal identification strategy. This framing is appropriate for the intended use case: competition authorities use screening markers to prioritize investigations, not as conclusive evidence of wrongdoing.

A practical concern is that cartels may \textit{adapt} to the FL screen once it becomes known. If procurement authorities begin flagging firms with zero win rates and high participation counts, cartels may respond by allowing cover bidders to occasionally win small tenders, rotating cover-bidding duties across a larger pool of firms, or using cover bidders with some legitimate business activity. Such strategic adaptation is a general challenge for any screening methodology \citep{harrington2008detecting}. However, adaptation has costs: allowing cover bidders to win occasionally defeats the purpose of the cartel's bid allocation, and expanding the pool of cover bidders increases the risk of defection and detection through leniency programs. The FL screen is therefore most valuable as part of a multi-layered detection strategy, where different screens target different aspects of cartel behavior and adaptation to one screen is constrained by exposure to others.

Our findings suggest a practical \textbf{two-stage detection workflow} for competition authorities: (1) apply the FL screen to identify firms with suspicious participation-without-winning patterns, requiring only participation and outcome data; (2) for tenders flagged in stage 1, apply bid-level screens \citep{imhof2019detecting, chassang2022robust} to assess whether bid distributions exhibit anomalies consistent with collusion. The FL screen serves as a computationally inexpensive filter that reduces the volume of tenders requiring detailed bid-level analysis. In procurement systems where bid-level data are incomplete or unreliable---common in developing countries---the FL screen may be the only feasible quantitative tool.

Our findings have broader implications for competition policy. The FL screen is operationally simple: it requires only firm-level participation and outcome data, not detailed bid values. It can be implemented with minimal computational resources and does not require econometric expertise. The screen is particularly valuable in developing countries where bid-level data may be incomplete or unreliable, institutional capacity for sophisticated econometric analysis is limited, and minimum bidder requirements create strong incentives for cover bidding arrangements.

Several extensions merit future research. First, the FL screen should be validated in other procurement platforms and countries to assess its external validity. Our results are specific to BEC's institutional environment; the prevalence and form of cover bidding may differ in platforms with different rules (e.g., no minimum bidder requirements, stronger audit mechanisms, different bid visibility rules). Second, \textbf{network analysis} of co-bidding patterns among FL firms and their associated winners could reveal the structure of cartel organizations---identifying not just individual cover bidders but the complete bidding rings to which they belong. Preliminary analysis of co-bidding networks in our data suggests clustered participation patterns among FL firms, but a rigorous network analysis is beyond the scope of this paper. Third, the interaction between FL screening and leniency programs deserves theoretical and empirical investigation: does the availability of a screening tool that identifies cover bidders increase the incentive for cover bidders to defect through leniency programs? Fourth, the two-regime framework could be tested in settings with richer bid-level data, potentially distinguishing firms that serve as cover bidders in some tenders but compete genuinely in others.
